% Canonical native-TeX manuscript; compile from the repository root.
% Publication layout changes are reversible in the import manifest.
% Original Subset Sum paper style, recovered without changing its typography.
% Original main.tex SHA-256:
% ff77ed5a42d72a79197b6881a35cf7982457667b21d88bf5a1eeaea61689945a
% Original subset-sum.pdf SHA-256:
% bdeb95400d450f64cfc27ebbae255718b9a6a26607d0c1375091a4c7b3a4a431
\documentclass[11pt,a4paper]{article}
\usepackage[T1]{fontenc}
\usepackage[utf8]{inputenc}
\usepackage{lmodern}
\usepackage[a4paper,margin=25mm]{geometry}
\usepackage{amsmath,amssymb,amsthm}
\usepackage{booktabs,longtable,array}
\usepackage{microtype}
\usepackage[numbers,sort&compress]{natbib}
\usepackage[unicode,colorlinks=true,linkcolor=blue,citecolor=blue,urlcolor=blue,hypertexnames=false]{hyperref}
\urlstyle{same}
\setlength{\emergencystretch}{2em}
\setlength{\parskip}{3pt}
\clubpenalty=10000
\widowpenalty=10000
\displaywidowpenalty=10000
\allowdisplaybreaks

\usepackage{mathtools}
\newtheorem{theorem}{Theorem}[section]
\newtheorem{lemma}[theorem]{Lemma}
\newtheorem{corollary}[theorem]{Corollary}
\newtheorem{definition}[theorem]{Definition}
\newcommand{\F}{\mathbb F}
\newcommand{\SV}{\operatorname{SV}}
\newcommand{\wt}{\operatorname{wt}}
\newcommand{\pdeg}{\operatorname{pdeg}}
\newcommand{\Coeff}{\operatorname{Coeff}}
\newcommand{\Span}{\operatorname{span}}
\newcommand{\rank}{\operatorname{rank}}
\newcommand{\Mat}{\operatorname{Mat}}
\newcommand{\GL}{\operatorname{GL}}
\newcommand{\rowspan}{\operatorname{rowspan}}
\title{Deterministic Identity Testing for Constant-Width\\
Affine ROABP Orbits in Every Characteristic}
\author{icekylinx \and GPT-6 Astra Pro}
\date{27 October 2026}
\begin{document}\maketitle

\begin{abstract}
We give a self-contained candidate proof of a polynomial-size black-box
hitting set, for every fixed width, for the unknown invertible affine
orbit of read-once oblivious algebraic branching programs in arbitrary
characteristic. The local univariate matrices may have arbitrary generic
ranks and arbitrary degrees; only the scalar total degree is bounded.
The proof extends a coefficient-span argument to finite monomial Artin
rings. A box quotient then bounds the support of every monomial of a
lowest homogeneous component whenever all interval ranks are preserved.
The standard characteristic-order criterion for Hasse pivots bounds
every exponent of each such leading monomial. Conditioning on the other
blocks supplies the required filtered interface in an exterior-power
rank descent. The resulting explicit map is a moment curve plus
$K(w)=2^{O(w^2)}$ Shpilka--Volkovich pairs. Its basic grid has at most
$(nD+1)^{2K(w)+1}$ points for scalar total degree at most $D$.
\end{abstract}

\noindent\textbf{Status and attribution.}
The constituent notes and this self-contained integrated manuscript have
separate model-assisted internal reconstructions and exact-source audits,
including an additional adversarial review of the Artin reset/source step.
The mathematical source is frozen after those internal reviews. It remains
an internal candidate proof, not external peer review or a
publication-priority claim.
The generic Hasse-order closure used below is a standard $p$-adic
criterion, recorded for example in Garcia--Lax~\cite[Proposition 1.1]{GL92}
and in the St\"ohr--Voloch order-sequence context~\cite{SV86}; no novelty
is claimed for that criterion. We prove the precise rational-vector
version needed here rather than importing extra geometric hypotheses.

\section{Statement and oracle model}

Let $n,w,D\ge1$. Suppose that the nonzero polynomial presented to the
oracle has the form
\begin{equation}\label{eq:model}
 f(X)=uA_1(\ell_1(X)+c_1)\cdots A_n(\ell_n(X)+c_n)v,
 \qquad \deg f\le D.
\end{equation}
The matrices $A_i$ are $w$-by-$w$ univariate polynomial matrices over
a field $\F$, and $u,v$ are constant endpoints. Their generic ranks
and entry degrees are unrestricted. The forms $\ell_1,\ldots,\ell_n$
are the coordinates of an invertible linear transformation. The
representation, its order, and the affine transformation are unknown.
A representation of smaller width can be padded with zeros.

This is the affine-orbit model for one unknown ROABP order, as studied
by Bhargava and Ghosh~\cite{BG}; it does not assume a representation
in every variable order. The degree promise is on the scalar polynomial.
If local entry degrees, or latent scalar individual degrees, are bounded
by $\delta$, one may take $D=n\delta$. No truncation of the promised
representation is performed.

For $w\ge2$ and $1\le r<w$ define
\begin{align}
 W_{w,r}&=\binom wr,\label{eq:W}\\
 h_{w,r}&=(r+2)\prod_{s=r+2}^{w}
                    \left(\binom{s-1}{r-1}+1\right)-2,\label{eq:h}\\
 k_{w,r}&=(W_{w,r}^2-1)(2h_{w,r}+2),\qquad
 K(w)=\sum_{r=1}^{w-1}k_{w,r}.\label{eq:K}
\end{align}
An empty product is one, and $K(1)=0$. These constants are independent
of the characteristic and of every degree parameter. The construction gives
\[
 \begin{array}{c|rrrrrr}
 w&1&2&3&4&5&6\\ \hline
 K(w)&0&12&128&1500&23718&553854
 \end{array}
\]

Choose $n$ distinct nodes $a_i$ in a working field and their Lagrange
polynomials $\lambda_i(y)$, of degree at most $n-1$. Write
\[
 \SV_{k,i}(z,y)=\sum_{j=1}^{k}z_j\lambda_i(y_j),\qquad
 \gamma_i(t)=t^i.
\]

\begin{theorem}\label{thm:main}
In every characteristic, every nonzero polynomial \eqref{eq:model}
satisfies
\begin{equation}\label{eq:map}
 f\bigl(\gamma(t)+\SV_{K(w)}(z,y)\bigr)\ne0.
\end{equation}
The pure map $\SV_{K(w)+1}$ also hits the class. A Cartesian grid
for \eqref{eq:map} uses at most
\begin{equation}\label{eq:grid}
 (nD+1)(D+1)^{K(w)}\bigl((n-1)D+1\bigr)^{K(w)}
 \le(nD+1)^{2K(w)+1}
\end{equation}
oracle evaluations. A working field of size at least $nD+1$ supplies
all grids and nodes; extension-field oracle evaluations are permitted
when needed. Moreover
\[
 K(w)\le(w-1)2^{w(w+1)/2+2w}=2^{O(w^2)}\qquad(w\ge2).
\]
Thus for every fixed width the hitting set is polynomial in $n,D$.
Constants, including the cases $n=0$ or $D=0$, need one query.
\end{theorem}

This is an XP-type query bound in the width parameter: the polynomial
exponent depends on $w$. It does not establish an FPT bound
$g(w)(nD)^{O(1)}$ with a uniform exponent. The parameter dependence is
large and is not asserted to be optimal.
The theorem concerns invertible affine transformations, not singular
projections or orbit closures. All final degree counts refer to the
original scalar $f$, even though auxiliary minors may have much larger
degrees.

\section{Coefficient spaces in monomial Artin rings}

Let $E$ be a field. For disjoint variable blocks $X_i$, let
$\mathcal R_i=E[X_i]/I_i$, where $I_i$ is a proper monomial ideal
containing a positive power of every variable. A trivial block
$\mathcal R_i=E$ is also allowed. The surviving monomials form an
$E$-basis, and the augmentation ideal is nilpotent. Work in the tensor
product $\mathcal R=\bigotimes_i\mathcal R_i$ over $E$.
Support counts the blocks containing a nonconstant monomial.

For a row-valued or matrix-valued polynomial over these rings,
$\Coeff$ means the $E$-linear span of its actual coefficient rows
in the surviving monomial bases.

\begin{lemma}\label{lem:coeff}
Scalar multiplication, including multiplication by a polynomial using
future variable blocks, only takes $E$-linear combinations of existing
coefficient rows. The same holds for a compatible matrix multiplying
a row family on the left. Left multiplication by a square matrix with
invertible constant term preserves the coefficient-row span in both
directions.
\end{lemma}
\begin{proof}
Coefficient convolution expresses each resulting row as a linear
combination of original rows. Monomial truncation discards terms and
never identifies a higher monomial with a lower one. A matrix with
invertible constant term has an inverse in the finite local ring.
Apply the inclusion to the matrix and its inverse.
\end{proof}

\begin{definition}\label{def:unitrank}
An interval matrix has \emph{unit rank} $s$ if its constant matrix
has field rank $s$ and all its $(s+1)$-minors vanish in the tensor
ring. At $s=0$ this means that the matrix is zero; at full width the
condition on larger minors is empty.
\end{definition}

The constant matrix supplies an $s$-minor with nonzero constant
determinant, which is a unit. Thus the definition imposes the precise
algebraic rank identities needed below without defining a fraction-field
rank over a ring with zero divisors.

\begin{lemma}\label{lem:graph}
Let $P$ have unit rank $s$, and suppose an $s$-by-$s$ block $E_0$
has invertible constant term. After arranging it in the upper left,
write $P=\left[\begin{smallmatrix}E_0&B\\C&D\end{smallmatrix}\right]$.
Then $D=CE_0^{-1}B$. Consequently every row in the row module of $P$
is determined by its first $s$ coordinates; a zero active part forces
the entire row to be zero. Constant invertible cutwise gauges preserve
all unit-rank conditions.
\end{lemma}
\begin{proof}
Each entry of the Schur complement $D-CE_0^{-1}B$ is, up to sign,
an $(s+1)$-minor divided by the unit $\det E_0$, and hence vanishes.
The row-graph statement follows from
$P=\left[\begin{smallmatrix}I\\CE_0^{-1}\end{smallmatrix}\right]
E_0[I,E_0^{-1}B]$.
Cauchy--Binet preserves the vanishing larger minors under constant
invertible changes, and the constant rank is unchanged.
\end{proof}

No integral-domain assumption is present in these lemmas.

\section{Weighted concentration with a rank palette}

Fix a width $W$ and a set of allowed ranks
\[
 \mathcal P=\{0=p_0<p_1<\cdots<p_m=W\}.
\]
Assume every contiguous interval of a word
$A_1(X_1)\cdots A_n(X_n)$ over the tensor rings has unit rank equal
to its constant rank, and that rank belongs to $\mathcal P$.
Unused ranks may be included in the palette.

For $a=p_j\in\mathcal P$ put
\begin{equation}\label{eq:G}
 G(a)+1=\prod_{\ell=j+1}^{m}(p_\ell-p_{\ell-1}+1).
\end{equation}
In particular $G(W)=0$. For $b>a$ in the palette,
\begin{equation}\label{eq:recurrence}
 (b-a+1)(G(b)+1)\le G(a)+1.
\end{equation}
Indeed inserting an intermediate level splits a jump into positive
$u,v$ and replaces $u+v+1$ by $(u+1)(v+1)\ge u+v+1$.

Suppose constant cutwise gauges have made
\[
 A_i(0)=\operatorname{diag}(I_a,Q_i),\qquad a\in\mathcal P.
\]
Let the initial row space be
\[
 S_0=\rowspan(B_0)\oplus(0\times K_0),
\]
where the active part $H$ of $B_0$ has independent rows and spans a
space $T$. Assign cost zero to rows of $B_0$ and cost one to rows
$(0,p)$ with $p\in K_0$. A weighted budget $b$ allows coefficients
of $B_0A_1\cdots A_n$ supported on at most $b$ blocks and coefficients
from the initial passive rows supported on at most $b-1$ blocks.
Assume that the active projection of the full reachable coefficient-row
space $S_i$ remains exactly $T$ at every cut.

\begin{lemma}[Weighted stable phase]\label{lem:weighted}
Under these assumptions, weighted budget $G(a)$ spans $S_n$.
\end{lemma}
\begin{proof}
Induct on the number of palette levels above $a$. At $a=W$ there
are no passive coordinates. The unchanged active space is already
spanned by the constant-prefix rows, and $G(W)=0$ suffices.
We give the exact kernel and reset identities before the induction step.

\paragraph{Kernel dynamics.}
Propagate $B_i=B_{i-1}A_i(0)$; its active part remains $H$. Write
$S_i=\rowspan(B_i)\oplus(0\times K_i)$ and put $k=W-a$.
For $d=\dim T>0$, choose a field right inverse $J$ of $H$.
The active block of each local matrix is a unit, so write
\[
 A_i=\begin{bmatrix}
 E_i&E_iC_i\\D_iE_i&D_iE_iC_i+F_i
 \end{bmatrix},\qquad
 E_i(0)=I_a,\quad C_i(0)=D_i(0)=0,\quad F_i(0)=Q_i.
\]
Active stability gives
\[
 (H+B_{i-1,p}D_i)E_i=R_iH,\qquad R_i(0)=I_d,
\]
where the subscript $p$ denotes the passive block. Define
\begin{align}
 W_i&=R_i^{-1}B_{i-1}A_i=[H,B_{i,p}+Z_i],\label{eq:Wi}\\
 \widetilde F_i&=(I_k-D_iE_iJR_i^{-1}B_{i-1,p})F_i.\label{eq:Fi}
\end{align}
These functions use only the current variable block: $B_{i-1},H,J$
are field-constant coefficient lifts. Moreover $Z_i(0)=0$ and
$\widetilde F_i(0)=Q_i$. Lemma~\ref{lem:coeff} gives
$\Coeff(W_i)=\Coeff(B_{i-1}A_i)$. Each coefficient row of $Z_i$
is consequently a combination of actual rows of weighted cost at most
one.

For $p\in K_{i-1}$, stability implies $pD_iE_iJH=pD_iE_i$.
Direct subtraction gives
\[
 (0,p)A_i-(pD_iE_iJ)W_i=(0,p\widetilde F_i).
\]
Taking coefficient spans in both directions, and subtracting active
parts using the available constant rows $B_i$, yields
\begin{equation}\label{eq:KD}
 K_i=\Coeff(Z_i)+\Coeff(K_{i-1}\widetilde F_i).
\end{equation}
For a reachable state $a_{i-1}B_{i-1}+(0,p_{i-1})$, its coordinates
evolve by the exact identities
\begin{equation}\label{eq:EV}
 a_i=a_{i-1}R_i+p_{i-1}D_iE_iJ,\qquad
 p_i=p_{i-1}\widetilde F_i+a_iZ_i.
\end{equation}
If $d=0$, the $B,W,Z$ families are empty; active stability forces
reachable active outputs to zero, and the formulas reduce to
$K_i=\Coeff(K_{i-1}F_i)$. No right inverse is required in that case.

\paragraph{Reset across a zero passive interval.}
Suppose $Q_c\cdots Q_e=0$. The original interval
$P=A_c\cdots A_e$ has unit rank $a$, and its active pivot is a unit.
It is a graph over the active coordinates by Lemma~\ref{lem:graph}.
For fixed $p\in K_{c-1}$, the active part of $B_{c-1}P$ is
$R_{[c,e]}H$ with $R_{[c,e]}(0)=I$, and that of $(0,p)P$ is
coefficientwise in $T$. A regular initial active control $a_{c-1}$
therefore makes
\[
 (a_{c-1}B_{c-1}+(0,p))P
\]
have zero final active part. The graph identity makes its entire
output zero. Since $H$ has a constant right inverse, its final active
coordinate $a_e$ is zero as well. Unrolling \eqref{eq:EV} gives
\begin{equation}\label{eq:reset}
 p\widetilde F_c\cdots\widetilde F_e
 =-\sum_{j=c}^{e-1}a_jZ_j\widetilde F_{j+1}\cdots\widetilde F_e.
\end{equation}
The $j=e$ term is absent because $a_e=0$.
When $d=\dim T=0$, the reachable row already has zero active part
and the same reset statement holds without controls or sources.

The initial control may use future interval variables. This does not
violate reachability: by Lemma~\ref{lem:coeff}, multiplying an actual
prefix row by such a scalar only combines its coefficient rows.
Every controlled intermediate state remains coefficientwise in $S_i$,
and its passive remainder remains in $K_i$. The reset identity is used
only for an \emph{unweighted coefficient-span inclusion}. No support
cost is assigned to its possibly complicated scalar multipliers.

\paragraph{Positive passive constant rank.}
Suppose $q=\rank(Q_1\cdots Q_n)>0$, and put $b=a+q$.
This is the rank of the full constant interval, so $b\in\mathcal P$.
Constant gauges split a persistent $q$-dimensional passive core.
For explicit construction, choose $U_0$ complementary to the kernel
of $Q_1\cdots Q_n$, set $U_i=U_0Q_1\cdots Q_i$, choose a
complement $V_n$ of $U_n$, and put
\[
 V_i=\{v:vQ_{i+1}\cdots Q_n\in V_n\}.
\]
Then $U_i,V_i$ are complementary of dimensions $q,k-q$; each constant
map preserves this splitting, and maps the $U_i$ isomorphically.
The constants become $\operatorname{diag}(I_a,I_q,Q_i')$.

Let $\beta$ be the dimension of the initial passive kernel's projection
onto this new core. The enlarged reachable active space starts at
dimension $\dim T+\beta$, is nested, and has dimension at most
$\dim T+q$. Hence there are at most $q-\beta$ strict growth steps.
Every stable subphase now has persistent core dimension $b$, so
induction applies with budget $G(b)$.

Here is the support accounting. If active lifts cost at most $L$ and
passive-kernel generators at most $L+1$, a stable phase costs at most
$L+G(b)$. The next growth step adds one fresh block; a new active
basis and its kernel can both be selected at cost $L+G(b)+1$.
Substituting these coefficient witnesses is legitimate because later
blocks are disjoint from earlier ones. The tensor quotient changes
no cross-block coefficient convolution.

If $\beta=0$, the initial costs are zero and one. At most $q$ growth
steps give
\[
 (q+1)G(b)+q=(q+1)(G(b)+1)-1\le G(a)
\]
by \eqref{eq:recurrence}. If $\beta\ge1$, choose additional active
lifts from the initial kernel at cost one. The remaining kernel also
costs one, within the permitted cost two. At most $q-\beta$ growths
give
\[
 1+(q-\beta+1)G(b)+(q-\beta)
 =(q-\beta+1)(G(b)+1)
 \le q(G(b)+1)\le G(a).
\]
The last inequality uses $G(b)\ge0$ and \eqref{eq:recurrence}.
Only larger active levels $b>a$ have been used.

\paragraph{Zero passive constant rank.}
Suppose $Q_1\cdots Q_n=0$. Choose $c$ maximal with
$Q_c\cdots Q_n=0$. Every proper suffix $Q_{j+1}\cdots Q_n$,
$j\ge c$, has positive rank, including the empty suffix of rank $k$.
Equations \eqref{eq:KD} and \eqref{eq:reset} show that $K_n$ is
spanned by coefficients of
\[
 Z_j\widetilde F_{j+1}\cdots\widetilde F_n,\qquad j\ge c.
\]
For a coefficient row $p$ of $Z_j$, start the actual suffix
$A_{j+1}\cdots A_n$ from $\Span(B_j,(0,p))$, assigning weights
zero and one. Its old active projection remains $T$: it is contained
in the global reachable space, and its constant-prefix rows already
supply $T$. Its positive-passive-rank case, just proved, has budget
$G(a)$. Its kernel recurrence uses the same $B_i,W_i,Z_i,
\widetilde F_i$ and contains
$p\widetilde F_{j+1}\cdots\widetilde F_n$.

That response is therefore spanned by actual suffix coefficients
from $B_j$ of support at most $G(a)$, or from $p$ of support at most
$G(a)-1$. In the original word $B_j$ costs zero and $p$ costs at most
one. Its source cost is already paid by the weight one, so the total
budget is $G(a)$. This is not circular induction at level $a$:
the positive-rank case invoked only levels $b>a$. Adding the constant
rows $B_n$ completes the span of $S_n$.

\paragraph{Preservation of the palette.}
Every subphase and source restart uses a contiguous subword of the
same full matrices. It changes an initial row space, not the matrices.
The cutwise gauges are invertible field-constant changes and preserve
each interval rank and its unit-pivot identities. Each newly enlarged
core dimension is the rank of an actual full constant interval, hence
belongs to $\mathcal P$. The corrected passive operators
$\widetilde F_i$ are used only in the exact algebra; no rank palette
is assumed for their products. These observations justify all
induction calls and finish the proof.
\end{proof}

\begin{theorem}[Palette row concentration]\label{thm:row}
For every fixed row $u$, the coefficient-row span of
$uA_1\cdots A_n$ is generated by actual coefficients supported on at
most
\begin{equation}\label{eq:paletteh}
 h(\mathcal P)=\prod_{j=1}^{m}(p_j-p_{j-1}+1)-2
\end{equation}
blocks.
\end{theorem}
\begin{proof}
Use Lemma~\ref{lem:weighted} with empty active core $a=0$ and initial
passive space $\Span(u)$ of weight one. Weighted budget $G(0)$ means
ordinary support budget $G(0)-1$, which is \eqref{eq:paletteh}.
A full interval of unit rank zero is the zero matrix and causes no
exception.
\end{proof}

For the full palette $\{0,1,\ldots,W\}$, this gives $2^W-2$.
The palette refinement will be substantially smaller on compound words.

\section{Box quotients bound every leading support}

Return to polynomial or regular rational univariate matrices
$A_i(Y_i)$ over a field $E$. Suppose every contiguous interval has
generic rank equal to its constant rank at zero, and all these ranks
belong to $\mathcal P$. Let $u,v$ be any fixed vectors over $E$ and
assume $g(Y)=uA_1(Y_1)\cdots A_n(Y_n)v\ne0$.

\begin{theorem}\label{thm:leading-support}
Every nonzero monomial in the first ordinary-homogeneous component
of $g$ has support at most $h(\mathcal P)$.
\end{theorem}
\begin{proof}
A nonzero constant component is immediate. Otherwise let $m>0$ be
the first degree and choose any nonzero coefficient $g_\alpha$ with
$|\alpha|=m$. Set variables outside its support to zero and pass to
the tensor box
\[
 \bigotimes_{i:\alpha_i>0} E[Y_i]/(Y_i^{\alpha_i+1}).
\]
Every surviving monomial is a divisor of $Y^\alpha$. Each proper
divisor has total degree below $m$ and zero scalar coefficient. Thus
the entire scalar in this quotient is exactly $g_\alpha Y^\alpha$,
which is nonzero.

All interval hypotheses survive the quotient. Larger minors vanish
identically before quotienting, while every chosen constant nonzero
pivot stays a unit. Regular rational denominators are units as well.
Theorem~\ref{thm:row} therefore applies in this particular box. If
$|\operatorname{supp}\alpha|>h(\mathcal P)$, every scalar coefficient
permitted by its row-span bound is zero, forcing the entire scalar to
be zero. This contradicts $g_\alpha\ne0$.
\end{proof}

The box step is essential. Ordinary unfiltered coefficient concentration
alone would not bound every leading support: it could use a low-support
coefficient of a much higher total degree. The quotient removes those
coefficients while retaining all needed Schur identities.
The endpoints here are arbitrary constants over $E$; they may depend
on other shifts or parameters, as long as they are fixed in the residual
variables being considered.

\section{Standard Hasse orders and the filtered exponent bound}

In characteristic $p>0$, let $\wt_p(a)$ be the sum of the base-$p$
digits of $a$. The digit degree of a monomial is the sum of the digit
weights of its exponents. In characteristic zero use $\wt_0(a)=a$.
For a rational function, Hasse derivatives are defined by
\[
 F(T+Z)=\sum_{j\ge0}D_T^{[j]}F(T)Z^j.
\]
Their product and composition rules hold over every field.

\begin{lemma}[Generic pivot closure]\label{lem:pivots}
Let $F(T)$ be a nonzero row vector of $q$ rational functions over
an arbitrary original field $F_0$. Call $\nu$ a pivot if
$D^{[\nu]}F$ does not lie in the $F_0(T)$-span of all earlier jets.
There are at most $q$ pivots. If $\nu$ is a pivot and
$\binom\nu\mu\ne0$ in $F_0$, with $0\le\mu\le\nu$, then
$\mu$ is a pivot. Consequently every pivot satisfies
\[
 \wt_p(\nu)\le q-1
\]
in positive characteristic, and $\nu\le q-1$ in characteristic zero.
\end{lemma}
\begin{proof}
If $\mu$ were not a pivot, write
\[
 D^{[\mu]}F=\sum_{j<\mu}a_j(T)D^{[j]}F.
\]
Apply $D^{[\nu-\mu]}$. The left side is
$\binom\nu\mu D^{[\nu]}F$. The product rule on the right gives
only jets of indices $j+\nu-\mu-e<\nu$, for $e\ge0$.
This contradicts the pivot property of $\nu$.

In characteristic $p$, the nonzero binomial coefficients are exactly
the digit subindices, as follows from the Frobenius expansion of
$(1+Z)^\nu$. If $\nu=\sum_e\nu_ep^e$, all
$\prod_e(\nu_e+1)$ digit subindices are pivots. Thus
\[
 \wt_p(\nu)+1\le\prod_e(\nu_e+1)\le q.
\]
In characteristic zero every index $0,\ldots,\nu$ is a pivot.
This is the standard $p$-adic order criterion cited above, proved here
for the precise rational-vector setting.
\end{proof}

\begin{lemma}\label{lem:moving-functional}
Let $C\in E$ be transcendental over $F_0$, and let
$\Lambda:E^q\to E$ be any $E$-linear functional. If
$\Lambda(F(C+Y))$ is nonzero and first becomes nonzero at order $a$,
then $a$ is a generic pivot of $F$. In particular $\wt_p(a)\le q-1$,
with ordinary weight in characteristic zero.
\end{lemma}
\begin{proof}
Every original-field rational jet relation is defined at $C$ because
its nonzero denominator cannot vanish there. Every nonzero finite jet
minor remains nonzero as well, so all successive jet ranks and pivots
are preserved at $C$. If $a$ were not a pivot, its specialized jet
would be a linear combination of earlier ones, all annihilated by
$\Lambda$, contradicting its nonzero value.
\end{proof}

The functional may depend on $C$ and other correlated parameters.
Neither perfection of $F_0$ nor transcendence of $C$ over the enlarged
field $E$ is assumed.

\begin{lemma}[Every leading exponent]\label{lem:exponents}
Let $A_i(T_i)$ be rational univariate matrices of width $W$ over
$F_0$, and let every $C_i\in E$ be transcendental over $F_0$.
For arbitrary fixed $E$-endpoints, suppose
\[
 g(Y)=uA_1(C_1+Y_1)\cdots A_n(C_n+Y_n)v\ne0.
\]
For every nonzero monomial $Y^\alpha$ of its first homogeneous
component, every coordinate satisfies
$\wt_p(\alpha_i)\le W^2-1$. Local ranks are unrestricted.
\end{lemma}
\begin{proof}
Let $m=|\alpha|$. Hold all exponents other than that of $Y_i$ at
their chosen values. The remaining scalar coefficient at order $a$
is an $E$-linear matrix-entry functional of $D^{[a]}A_i(C_i)$:
its left and right vectors are the products of the fixed outside
coefficient matrices and endpoints. At every $a<\alpha_i$ this
coefficient is zero, since its full total degree is below $m$.
At $a=\alpha_i$ it is nonzero. Apply
Lemma~\ref{lem:moving-functional} to the vector of the $W^2$
original matrix entries. Constants and empty variable blocks have
only order-zero coefficients and cause no exception.
\end{proof}

This is an ordinary-order filtration argument for every leading
monomial. It is stronger than selecting a low-digit-weight coefficient
somewhere in the expansion. Only individual original-field
transcendence is needed; there is no derivative-separating or
common-translation premise in Lemma~\ref{lem:exponents}.

\section{Affine detection of a bounded leading component}

\begin{lemma}\label{lem:affine-detection}
Let a nonzero polynomial or regular series have a first homogeneous
component of digit degree at most $k$. After any invertible linear
residual change of coordinates, $\SV_k$ detects that component and
the whole series remains nonzero under the substitution.
\end{lemma}
\begin{proof}
An invertible linear change preserves nonzeroness. Digit degree cannot
increase: expand each exponent into its digit copies, use Frobenius
on every copy, and note that collisions can only decrease digit weight.
In characteristic zero this is ordinary total degree preservation.
Every resulting monomial has support at most $k$.

The image of $\SV_k$ contains every coordinate subspace of dimension
at most $k$: specialize its direction variables to the corresponding
interpolation nodes and unused amplitudes to zero. Restricting to
the support of a nonzero monomial retains that coefficient, so one of
these subspaces has a nonzero restriction. Thus the first component's
SV image is nonzero. Higher homogeneous components have larger total
degree in the fresh amplitudes and cannot cancel it. A nonzero constant
component needs no fresh variables.
\end{proof}

All statements remain valid over a coefficient extension generated by
earlier symbolic seeds; evaluation at the interpolation nodes still
gives the identity matrix for the Lagrange basis.

\section{The interval-rank descent}

Adjoin the constant width-$w$ factors
\[
 E_{\mathrm{in}}=e_1u,\qquad E_{\mathrm{out}}=ve_1^T
\]
to the original latent word. Its full product is
$f_{\mathrm{latent}}(T)E_{11}$. These are empty variable blocks and
add no observed coordinates. Every interval mentioned in the descent
belongs to this \emph{unshifted augmented word}.

At any stage let $G$ be the moment curve plus already introduced
SV pairs, with all seeds retained as independent indeterminates. Put
\[
 C_i=c_i+\ell_i(G).
\]
Every $C_i$ is transcendental over the original working field: its
$t$-part is a nonzero linear combination of $t,t^2,\ldots,t^n$,
whereas all additional SV terms are independent of $t$.
This remains true in every characteristic, even if an ordinary
derivative of that combination vanishes.

Before the rank-$r$ step the induction invariant is
\begin{equation}\label{eq:invariant}
 \begin{gathered}
 \text{every augmented interval of generic rank greater than }r\\
 \text{retains that rank at the current }C.
 \end{gathered}
\end{equation}
At $G=\gamma$, every rank-$w$ interval consists solely of full-rank
locals. Their nonzero original-field univariate determinants survive
at the individually transcendental $C_i$, so their product does too.
Constant invertible factors also survive. This starts the induction
at $r=w-1$. Previously preserved minors stay nonzero after adding
fresh pairs, because setting their amplitudes to zero recovers the
old shift.

\subsection{Greedy blocks inside one target interval}

Fix a contiguous interval $I$ of the unshifted augmented word with
generic rank exactly $r$. Every subinterval of $I$ has rank at least
$r$, since the rank of a full product is at most that of any contiguous
subproduct. Starting at the left, take a shortest prefix of rank $r$,
then repeat on the remaining suffix as long as such a prefix exists.
Write
\begin{equation}\label{eq:greedy}
 I=B_1\cdots B_qT,\qquad B_j=C_jA_{i_j},
\end{equation}
where $A_{i_j}$ is the last single factor of that shortest prefix.
There is at least one block. The chunks $C_j$ and leftover tail $T$
have generic rank greater than $r$, or are empty. Indeed a shortest
prefix cannot reach below $r$, and all its proper prefixes have
higher rank; stopping means the remaining full suffix has higher
rank. Every subinterval within one of these chunks has higher rank
as well.

The following words all have generic rank exactly $r$:
\begin{equation}\label{eq:tests}
 B_j,\qquad B_jB_{j+1}\ (j<q),\qquad B_qT.
\end{equation}
They are subintervals of $I$, giving rank at least $r$, and each
contains a rank-$r$ block, giving rank at most $r$. Choose for each
one a nonzero $r$-minor. The algorithm does not need to discover this
choice; it is used only in the proof.

\subsection{Exterior chunks with a preserved palette}

Put $W=\binom wr$ and replace each local matrix by its $r$-th
compound matrix. Cauchy--Binet gives
\[
 \bigwedge^r(AB)=(\bigwedge^rA)(\bigwedge^rB)
\]
in every characteristic. A matrix of rank $s$ has compound rank
$\binom sr$, by reducing it over its field to a diagonal rank-$s$
matrix. Thus each selected minor in \eqref{eq:tests} is a nonzero
scalar entry of a width-$W$ compound word whose locals are still
original-field univariate polynomial matrices.

By \eqref{eq:invariant}, every interval inside $C_j$ or $T$ has
preserved original rank $s>r$. Its compound interval has preserved
rank $\binom sr$. Hence every prepared compound chunk satisfies
Theorem~\ref{thm:leading-support} with palette
\[
 \mathcal P_{w,r}=\left\{0,\binom{r+1}{r},
                    \binom{r+2}{r},\ldots,\binom wr\right\}.
\]
Its first positive rank is $r+1$, and later gaps are
$\binom{s-1}{r-1}$. Formula \eqref{eq:paletteh} therefore gives
exactly $h(\mathcal P_{w,r})=h_{w,r}$ from \eqref{eq:h}.
An empty chunk is the identity and contributes support zero.

\subsection{Conditioning the leading monomial of a test}

Let $g(Y)$ be one selected nonzero compound-entry test after the
current shift $C$, in independent latent residual variables. It
remains a nonzero polynomial: translation by $C$ is an automorphism
over the old-seed field. If its constant term is nonzero, it is already
preserved. Otherwise choose \emph{any} nonzero monomial $Y^\alpha$
of its first degree $m$.

Fix the exponents outside one prepared chunk. Coefficient extraction
there leaves a scalar contraction of the chunk with a fixed row and
column over the old-seed field. These endpoints are products of the
fixed outside local coefficient matrices; they can depend on those
exponents and all shifts, but not on the chunk's residual variables.
The chosen chunk coefficient is nonzero. Every coefficient of smaller
\emph{total chunk degree} is zero, since it would yield a nonzero
coefficient of the full test below degree $m$ with the same outside
exponents. The chunk part of $\alpha$ therefore belongs to its first
homogeneous scalar component.

Theorem~\ref{thm:leading-support} bounds its support by $h_{w,r}$.
This is valid for arbitrary coefficient-field endpoints, and also
when the chunk degree is zero. Each word in \eqref{eq:tests} contains
at most two prepared chunks and at most two singleton layers:
\[
 C_jA_{i_j},\qquad
 C_jA_{i_j}C_{j+1}A_{i_{j+1}},\qquad
 C_qA_{i_q}T.
\]
Thus every leading monomial of every test has support at most
$2h_{w,r}+2$.

Apply Lemma~\ref{lem:exponents} independently to the whole compound
test. Its local matrices have width $W$, are over the original field,
and have individually transcendental shifts. Every leading exponent
has digit weight at most $W^2-1$, including those at singleton layers
where no rank-preservation premise has been used. Consequently the
\emph{entire} first homogeneous component of each test has digit
degree at most
\[
 (W^2-1)(2h_{w,r}+2)=k_{w,r}.
\]
No multivariate boundary vector has been factored. Both support and
exponent bounds are applied to first ordinary degrees, rather than
to unfiltered coefficient witnesses.

\subsection{One stage preserves rank \texorpdfstring{$r$}{r}}

Add a fresh $\SV_{k_{w,r}}$ in the observed coordinates. The latent
residual change is the original invertible linear transformation, so
Lemma~\ref{lem:affine-detection} preserves every selected test.
The same map works simultaneously for all tests in all target
intervals: they form a finite family of nonzero rational functions
in the common old and fresh seed field. No seed specialization is
made during the symbolic induction.

Every word in \eqref{eq:tests} now has specialized rank at least $r$
because its chosen minor survives, and at most $r$ because its larger
minors vanish identically. These ranks are therefore exactly $r$.
The following elementary chain fact completes the step.

\begin{lemma}\label{lem:chain}
If $B_1,\ldots,B_q$ have rank $r$, every adjacent product has rank
$r$, and $B_qT$ has rank $r$, then $B_1\cdots B_qT$ has rank $r$.
\end{lemma}
\begin{proof}
Inductively the prefix product through $B_j$ has rank $r$ and the
same row space as $B_j$. Multiplication by $B_{j+1}$ has rank $r$
because $B_jB_{j+1}$ does; the new row space equals that of
$B_{j+1}$. At the last block it therefore has the same row space
as $B_q$, whose multiplication by $T$ has rank $r$.
\end{proof}

The proof uses the tail test $B_qT$ separately, so there is no
uncounted three-chunk test after absorbing the tail. We have established
\eqref{eq:invariant} for the next lower rank while preserving all
previous ranks. Descend through $r=w-1,\ldots,1$. The final map is
the sum of the moment curve and exactly $K(w)$ fresh SV pairs.
The full augmented interval has generic rank one because $f\ne0$;
its rank preservation gives $f(G)\ne0$. If $w=1$, the initial curve
already preserves the product of nonzero univariate scalars.

\section{Quantitative bounds and the pure map}

For $1\le r<w$, we have $r+2\le2^{r+1}$ and
$\binom{s-1}{r-1}+1\le2^s$. Thus
\[
 h_{w,r}+2\le2^{\sum_{s=r+1}^{w}s}
             \le2^{w(w+1)/2-1}.
\]
Since $W_{w,r}\le2^w$, the safe estimate
\[
 K(w)\le(w-1)2^{w(w+1)/2+2w}
\]
follows. Much smaller constants are possible in several small-width
special cases; no such refinements are dependencies here.

The observed polynomial after the final substitution has individual
seed degrees at most $nD$ in $t$, $D$ in each amplitude, and
$(n-1)D$ in each direction. Cartesian interpolation proves
\eqref{eq:grid}. It evaluates only the original scalar polynomial on
an explicit polynomial map. Numerical points need not avoid the
auxiliary pivots used in the proof, and the degrees of auxiliary
minors do not enter this grid calculation.

For the pure map $\SV_{K(w)+1}$, specialize its first pair to
$z_0=t,y_0=t$ in every characteristic. A nonzero latent combination
of this first pair is $c+tq(t)$, where $q$ is a nonzero combination
of the independent Lagrange polynomials. It is nonconstant, hence
transcendental over the original field. All later SV stages are
independent of $t$. Therefore the entire preceding proof applies to
this first-pair specialization, and the unspecialized pure SV image
is nonzero as well. A basic pure-map grid uses at most
\[
 \bigl((D+1)((n-1)D+1)\bigr)^{K(w)+1}
 \le(nD+1)^{2K(w)+2}
\]
points. This completes Theorem~\ref{thm:main}.

\section{What the proof does and does not use}

The rank profile is used through constant-unit pivots and vanishing
larger minors, so it survives monomial box quotients. The weighted
reset identity is a coefficient-span statement; it does not claim
preservation of a matrix image variety or a nonlinear cone. The
Hasse argument is applied after filtering by the first ordinary total
degree, so it bounds each exponent of every leading monomial. It is
not the invalid inference that one low-digit-weight coefficient must
survive an unknown linear change.

The rank descent always decomposes a particular target interval of
generic rank $r$. That is what prevents a greedy prefix from jumping
below $r$, and ensures every adjacent-block and tail test has rank
exactly $r$. Conditioning on the outside coefficients leaves constant
endpoints for the chosen chunk, even when those endpoints depend on
all old seeds. The coefficient field in the generic Hasse relations
remains the original field, where each shifted latent argument is
transcendental. These are the interfaces on which the proof depends.

The final conclusion is an internal candidate proof with polynomial
size for fixed width. It does not give a polynomial bound uniform in
unbounded width, assert an optimal parameter count, or cover singular
affine projections or orbit closures.

\begin{thebibliography}{9}\raggedright
\bibitem{BG}
Vishwas Bhargava and Sumanta Ghosh.
\emph{Improved Hitting Set for Orbit of ROABPs}.
Computational Complexity, 2022.
\href{https://doi.org/10.1007/s00037-022-00230-9}{doi:10.1007/s00037-022-00230-9}.
Primary preprint:
\url{https://eccc.weizmann.ac.il/report/2021/062/revision/2/download/}.

\bibitem{GL92}
Arnaldo Garcia and R.\ F.\ Lax.
\emph{Weierstrass Weight of Gorenstein Singularities with One or Two Branches}.
1992 preprint, Proposition 1.1.
\url{https://arxiv.org/pdf/alg-geom/9208003}.

\bibitem{SV86}
K.-O. St\"ohr and J.\ F.\ Voloch.
\emph{Weierstrass Points and Curves Over Finite Fields}.
Proceedings of the London Mathematical Society, 52(1):1--19, 1986.
\href{https://doi.org/10.1112/plms/s3-52.1.1}{doi:10.1112/plms/s3-52.1.1}.
\end{thebibliography}
\end{document}
